\oldpage[IV-4]{155}
étant, il résulte d'abord de (\textbf{0}, 16.1.5) et (\textbf{0}, 16.1.4.1) que la relation $\dim(\sh{O}_{X',x'})\geq 2$ entraîne $\dim(\sh{O}_{X,x})\geq 2$. Inversement, supposons que $\dim(\sh{O}_{X,x})\geq 2$ et que $A$ soit \emph{universellement caténaire} (5.6.2); alors il en est de même de $\sh{O}_{X,x}$ (5.6.3), donc $\dim(\sh{O}_{X'',x''})=\dim(\sh{O}_{X,x})\geq 2$ en vertu de (5.6.10) et on a donc alors $\dim(\sh{O}_{X',x'})\geq 2$. La même conclusion est valable si $\sh{O}_{X,x}$ est \emph{géométriquement unibranche}, car la fibre du morphisme $X'\to X$ au point $x$ est alors réduite à un seul point, et $\sh{O}_{X',x'}$ est une $\sh{O}_{X,x}$-algèbre \emph{entière}; il suffit donc d'appliquer (\textbf{0}, 16.1.5).
\end{remark}

\begin{lemma}[18.9.7.7]
\label{IV.18.9.7.7}
Soient $A$ un anneau local noethérien, de dimension $\geq 2$, intègre et géométriquement unibranche, $Y=\Spec(A)$, $y$ le point fermé de $Y$. Soient $X$ un préschéma localement noethérien et connexe, $f:X\to Y$ un morphisme plat. Alors, pour toute partie fermée $T\subset f^{-1}(y)$, $X\setmin T$ est connexe.
\end{lemma}

Le même raisonnement qu'au début de (18.9.7.2) ramène à prouver que, pour tout point $x\in f^{-1}(y)$, $\Spec(\sh{O}_{X,x})\setmin\{x\}$ est connexe. On peut donc se borner au cas où $X=\Spec(B)$, $B$ étant un anneau local noethérien et l'homomorphisme $A\to B$ correspondant à $f$ étant local; posons $U=Y\setmin\{y\}$, et notons que $V=f^{-1}(U)$ est dense dans $X$ (2.3.10); il suffit donc de prouver que $V$ est \emph{connexe}. Gardons les notations de (18.9.7.5), et posons donc $A^{(1)}=\Gamma(U,\sh{O}_Y)$, $Y_1=\Spec(A^{(1)})$, $X_1=X\times_Y Y_1=\Spec(B\otimes_A A^{(1)})$, de sorte qu'on a le diagramme commutatif
\[
  \xymatrix{
    X \ar[d]_{f} & X_1 \ar[l]_{g_1} \ar[d]^{f_1} \\
    Y & Y_1 \ar[l]^{g}
  }
\]

En vertu de (2.3.1) et de la définition de $A^{(1)}$, on a un isomorphisme canonique $\Gamma(V,\sh{O}_X)\xrightarrow{\sim}B\otimes_A A^{(1)}$, et il suffit donc de prouver que ce dernier anneau est \emph{local}, un tel anneau ne pouvant être le produit de deux anneaux non réduits à 0 (\textbf{III}, 7.8.6.1). Mais on a vu (18.9.7.5) que $A^{(1)}$ est une sous-$A$-algèbre de la clôture intégrale $A'$ de $A$, puisque $\dim(A)\geq 2$. L'hypothèse que $A$ est géométriquement unibranche entraîne que le morphisme $\Spec(A')\to\Spec(A)$ est radiciel au point $y$ (6.15.3); \emph{a fortiori} $g:Y_1\to Y$ est radiciel au point $y$, et par suite $g_1$ est radiciel au point fermé $x$ de $X$; comme en outre $g_1$ est un morphisme entier, les points fermés de $X_1$ sont au-dessus de $x$, donc il n'existe qu'un seul point fermé de $X_1$, ce qui achève la démonstration de (18.9.7.7).

Il est clair que (18.9.7.7) démontre (18.9.7.3) lorsque $\dim(A)\geq 2$ (même sans supposer les fibres $X_y$ géométriquement réduites).

C) \emph{Cas où $\dim(A)=1$.} On a vu au début de la preuve de (18.9.7.3) que l'on peut supposer $\dim(B\otimes_A k)=0$, et par platitude (6.1.1.1), on a donc
\[
  \dim(B)=\dim(A)+\dim(B\otimes_A k)=1.
\]
Puisque $A$ est intègre, $Y$ se compose de deux points, le point générique $\eta$ et le point fermé $a$. En outre, puisque $f$ est plat, toute composante irréductible de $X$ domine $Y$ (2.3.4), donc $X\setmin\{b\}$, qui est l'ensemble des points maximaux $\xi_i$ de $X$, est égal à
\oldpage[IV-4]{156}
l'ensemble sous-jacent à $f^{-1}(\eta)$, et la fibre $f^{-1}(\eta)$ est somme des $\Spec(\boldsymbol{k}(\xi_i))$. L'hypothèse sur les $X_y$ et le fait que ces fibres sont de dimension 0, montrent qu'elles sont \emph{géométriquement régulières} (6.7.6), et \emph{a fortiori} géométriquement normales, autrement dit $f$ est un \emph{morphisme normal}. Mais puisque $A$ est intègre et géométriquement unibranche, il résulte alors de (6.15.10) que $B$ est intègre et géométriquement unibranche, donc $f^{-1}(\eta)=X\setmin\{b\}$ est réduit à un seul point, et \emph{a fortiori} connexe; ceci termine donc la démonstration de (18.9.7.3) et celle de (18.9.7).

\begin{remark}[18.9.7.8]
\label{IV.18.9.7.8}
Dans le cas C) de la démonstration de (18.9.7.3), on peut éviter de faire appel au délicat résultat (6.15.10) en raisonnant de la façon suivante: puisque $A$ est intègre et de dimension 1, il résulte du théorème de Krull-Akizuki (Bourbaki, \emph{Alg. comm.}, chap.~VII, \textsection~2, n\textsuperscript{o}~5, prop.~5) que sa clôture intégrale $A'$ est un anneau \emph{noethérien}. Le même raisonnement qu'au début de la démonstration de (18.9.7) montre alors que $B'=B\otimes_A A'$ est un anneau local; en outre, par platitude, $B'$ est contenu dans l'anneau total des fractions $R$ de l'anneau réduit $B$ (3.3.5). Or, le raisonnement qui prouve le théorème de Krull-Akizuki (Bourbaki, \emph{loc. cit.}) s'applique également à un anneau local noethérien \emph{réduit} de dimension 1, et montre que pour un tel anneau $B$, tout anneau compris entre $B$ et son anneau total des fractions est \emph{noethérien}. L'anneau $B'$ étant un anneau local noethérien et le morphisme $\Spec(B')\to\Spec(A')$ étant normal, $B'$ est un anneau intègre et intégralement clos (6.5.4), et on conclut comme au début de la preuve de (18.9.7).
\end{remark}

\begin{corollary}[18.9.8]
\label{IV.18.9.8}
Avec les notations de \textup{(18.9.7)}, supposons vérifiée la condition \textup{(i)} de \textup{(18.9.7)} et \emph{l'une} des deux conditions suivantes:
\begin{enumerate}
  \item[\textup{(ii$'$)}] $A$ est un anneau strictement local \textup{(18.8.2)}.
  \item[\textup{(ii$''$)}] $A$ est un anneau local hensélien et $\boldsymbol{k}(B)$ est une extension primaire de $\boldsymbol{k}(A)$.
\end{enumerate}

Alors \emph{toutes} les fibres $X_y$ \textup{($y\in Y$)} sont géométriquement connexes.
\end{corollary}

Soit $\mathfrak p$ l'idéal premier de $A$ correspondant au point $y$, et posons $A'=A/\mathfrak p$ et $B'=B\otimes_A A'=B/\mathfrak pB$, qui sont évidemment des anneaux locaux noethériens; il est clair que le morphisme $f':\Spec(B')\to\Spec(A')$ vérifie la condition (i) de (18.9.7); comme $y$ est le point générique de $\Spec(A')$ et que $\boldsymbol{k}(A')=\boldsymbol{k}(A)$, $\boldsymbol{k}(B')=\boldsymbol{k}(B)$, on voit qu'il suffit de vérifier que, dans le cas (ii$'$) (resp. (ii$''$)), $A'$ est géométriquement unibranche (resp. unibranche). Mais $A'$ est intègre, et dans le cas (ii$'$) (resp. (ii$''$)) il est strictement local (resp. hensélien) en vertu de (18.5.10). On en conclut que $A'$ est géométriquement unibranche (resp. unibranche) en vertu de (18.8.15) (resp. (18.6.12)).

\begin{corollary}[18.9.9]
\label{IV.18.9.9}
Soit $A$ un anneau local noethérien hensélien, universellement japonais \textup{(i.e. ((7.6.4) et (7.7.2)), dont les \emph{fibres formelles} (7.3.13) sont géométriquement réduites)}; alors les fibres formelles de $A$ sont géométriquement connexes.
\end{corollary}

Il suffit d'appliquer (18.9.8) au cas où $B=\widehat A$, puisque $B$ est un $A$-module plat et que $\boldsymbol{k}(B)=\boldsymbol{k}(A)$.

\begin{remark}[18.9.10]
\label{IV.18.9.10}
Les démonstrations de (18.9.4) et (18.9.5) n'utilisent que le fait que les fibres formelles de $A$ sont géométriquement connexes. Les conclusions
\oldpage[IV-4]{157}
de ces deux propositions sont donc encore valables lorsqu'on suppose l'anneau local $A$ noethérien, hensélien et universellement japonais.
\end{remark}

\begin{corollary}[18.9.11]
\label{IV.18.9.11}
Soient $X$, $Y$ deux préschémas localement noethériens, $Y$ étant supposé géométriquement unibranche et $X$ connexe. Soit $f:X\to Y$ un morphisme réduit \textup{(i.e. (6.8.1) plat et à fibres géométriquement réduites)}. Alors, pour toute partie fermée $T$ de $X$ telle que $f(T)$ soit rare dans $Y$, $X\setmin T$ est connexe.
\end{corollary}

Comme $f$ est plat et que $f(T)$, étant rare, ne peut contenir de point maximal de $Y$, $T$ ne peut contenir de point maximal de $X$ (2.3.4), donc est rare dans $X$. Utilisant (15.5.6.1), il suffit de montrer que, pour tout $x\in T$, $\Spec(\sh{O}_{X,x})\setmin\{x\}$ est connexe. Posons $y=f(x)$, qui n'est pas un point maximal de $Y$. Comme l'anneau $\sh{O}_{Y,y}$ est géométriquement unibranche, on peut appliquer (18.9.7) au morphisme
\[
  f_1:\Spec(\sh{O}_{X,x})\longrightarrow\Spec(\sh{O}_{Y,y}),
\]
qui vérifie les conditions (i) et (ii) de ce théorème, et à la partie fermée $\{x\}$ de $\Spec(\sh{O}_{X,x})$, puisque $f_1(\{x\})=\{y\}$ est rare dans $\Spec(\sh{O}_{Y,y})$.

\subsection{Préschémas étales sur un préschéma géométriquement unibranche ou normal.}
\EGAReaderTarget{ega4.c04.subsec.18-10}{subsection:IV.18.10}

\begin{theorem}[18.10.1]
\label{IV.18.10.1}
Soient $f:X\to Y$ un morphisme localement de présentation finie, $x$ un point de $X$, $y=f(x)$. Supposons que $Y$ soit intègre et géométriquement unibranche au point $y$. Alors, pour que $f$ soit étale au point $x$, il faut et il suffit que les deux conditions suivantes soient vérifiées:
\begin{enumerate}
  \item[\textup{(i)}] $f$ est non ramifié au point $x$;
  \item[\textup{(ii)}] l'homomorphisme $\sh{O}_{Y,y}\to\sh{O}_{X,x}$ est injectif.
\end{enumerate}

De plus, s'il en est ainsi, $X$ est intègre et géométriquement unibranche au point $x$.

Enfin, si $\sh{O}_{Y,y}$ est noethérien, on peut, dans le critère précédent, remplacer la condition \textup{(ii)} par la condition

\textup{(ii \emph{bis})} On a $\dim\sh{O}_{Y,y}\leq\dim\sh{O}_{X,x}$.
\end{theorem}

Le fait que, si $f$ est étale, $X$ soit intègre et géométriquement unibranche au point $x$ est un cas particulier de (17.5.7). Il est clair d'autre part que les conditions (i) et (ii) sont nécessaires pour que $f$ soit étale au point $x$, puisque $\sh{O}_{X,x}$ est alors un $\sh{O}_{Y,y}$-module fidèlement plat (17.6.1). Prouvons que ces conditions sont suffisantes.

Posons $A=\sh{O}_{Y,y}$, $A'={}^{hs}\!A$ (18.8.7), $Y'=\Spec(A')$, $X'=X\times_Y Y'$, $f'=f_{(Y)}:X'\to Y'$. Soient $y'$ le point fermé de $Y'$, $x'$ un point de $X'$ au-dessus de $x$ et de $y'$; puisque le morphisme $g:Y'\to Y$ est \emph{plat} (18.8.8), tout revient à voir que $f'$ est étale au point $x'$ (17.7.1, (ii)). L'hypothèse (ii) entraîne que le morphisme $\Spec(\sh{O}_{X,x})\to\Spec(\sh{O}_{Y,y})$ est dominant (\textbf{I}, 1.2.7); comme $\sh{O}_{Y,y}$ est intègre, $\Spec(\sh{O}_{Y,y})$ a un seul point générique $y_1$, et il y a donc une générisation $x_1$ de $x$ dans $X$ au-dessus de $y_1$. Le morphisme projection $X'\to X$ étant plat, il existe une générisation $x'_1$ de $x'$ au-dessus de $x_1$ (2.3.4); posons $y'_1=f'(x'_1)$; on a par suite $g(y'_1)=y_1$. Or, l'hypothèse sur $y$ entraîne que $A'$ est \emph{intègre} (18.8.15), donc $Y'$ a un seul point générique $\eta$, nécessairement au-dessus de $y_1$ puisque $g$
\oldpage[IV-4]{158}
est plat (2.3.4); d'autre part, $\eta$ est aussi point générique de $g^{-1}(y_1)$ ($\mathbf{0}_{\mathrm I}$, 2.1.8), et on sait par ailleurs que les fibres de $g$ sont des espaces \emph{discrets} (18.8.12), donc $g^{-1}(y_1)=\{\eta\}$, et par suite $y'_1=\eta$. Le morphisme $\Spec(\sh{O}_{X',x'})\to\Spec(\sh{O}_{Y',y'})$ restriction de $f'$ est donc dominant, et puisque $\sh{O}_{Y',y'}=A'$ est intègre, l'homomorphisme correspondant $\sh{O}_{Y',y'}\to\sh{O}_{X',x'}$ est injectif (\textbf{I}, 1.2.7). D'autre part, $f'$ est non ramifié au point $x'$, (donc (18.8.3, \emph{d)}), il y a un voisinage $U$ de $x'$ dans $X'$ tel que $f'|U$ soit une \emph{immersion fermée} de $U$ dans $Y'$; \emph{a fortiori} l'homomorphisme $\sh{O}_{Y',y'}\to\sh{O}_{X',x'}$ est surjectif, et puisqu'il est injectif, il est \emph{bijectif}. Mais alors $\sh{O}_{X',x'}$ est un $\sh{O}_{Y',y'}$-module \emph{plat}, et $f'$ est donc étale au point $x'$ (17.6.1).

Enfin, lorsque $\sh{O}_{Y,y}$ est noethérien, on sait que si $f$ est étale au point $x$, on a $\dim\sh{O}_{Y,y}=\dim\sh{O}_{X,x}$ (17.6.4). Réciproquement, supposons vérifiées les conditions (i) et (ii \emph{bis}), et montrons qu'elles entraînent (ii). Soit $\mathfrak J$ le noyau de l'homomorphisme canonique $\sh{O}_{Y,y}\to\sh{O}_{X,x}$; restreignant au besoin $Y$ à un voisinage de $y$, on peut supposer que $\mathfrak J=\mathscr{J}_y$, où $\mathscr{J}$ est un Idéal de $\sh{O}_Y$; si $Y_1$ est le sous-préschéma fermé de $Y$ défini par $\mathscr{J}$, on peut, en restreignant encore au besoin $X$ et $Y$ à des voisinages de $x$ et $y$ respectivement, supposer que $f$ se factorise en $X\xrightarrow{f_1}Y_1\to Y$ (\textbf{I}, 6.5.1). Il est clair que $f_1$ est encore non ramifié au point $x$, donc (17.4.1) quasi-fini en ce point; la démonstration de (5.4.1, (i)) prouve alors que l'on a $\dim(\sh{O}_{X,x})\leq\dim(\sh{O}_{Y,y}/\mathfrak J)$. Mais si l'on avait $\mathfrak J\neq0$, on en conclurait que $\dim(\sh{O}_{Y,y}/\mathfrak J)<\dim(\sh{O}_{Y,y})$ puisque $Y$ est intègre (\textbf{0}, 16.1.2.2); on aurait par suite $\dim\sh{O}_{X,x}<\dim\sh{O}_{Y,y}$ contrairement à l'hypothèse, ce qui prouve (ii).

\begin{remarks}[18.10.2]
\label{IV.18.10.2}
\begin{enumerate}
  \item[(i)] Lorsque $A=\sh{O}_{Y,y}$ est noethérien, et que l'on sait que son complété $\widehat A$ est intègre (par exemple si $A$ est régulier), on peut, dans la démonstration précédente, remplacer $A'$ par $\widehat A$.
  \item[(ii)] Lorsque $Y$ est localement noethérien, on peut donner de (18.10.1) une démonstration plus rapide, n'utilisant pas la hensélisation stricte, mais faisant intervenir les résultats assez délicats des \textsection\textsection~14 et 15. Comme $f$ est quasi-fini au point $x$ par hypothèse (17.4.1), on peut, en remplaçant $X$ par un voisinage ouvert de $x$, supposer que $f^{-1}(y)=\{x\}$. L'hypothèse (ii) entraîne, comme on l'a vu, l'existence d'une composante irréductible $Z$ de $X$ contenant $x$ et dominant l'unique composante irréductible $Y_0$ de $Y$ contenant $y$. Le morphisme $f|Z$ étant quasi-fini au point $x$, donc équidimensionnel en ce point (13.2.2), il résulte de l'hypothèse sur $Y$ et du critère de Chevalley (14.4.4) que $g=f|Z$ est universellement ouvert au point $x$. Par ailleurs puisque $f$ (donc aussi $g$) est non ramifié au point $x$, $g^{-1}(y)$ est géométriquement réduit sur $\boldsymbol{k}(y)$ (17.4.1); comme $\sh{O}_{Y,y}$ est intègre, on conclut de (15.2.3) que $g$ est \emph{plat} au point $x$, donc $f$ est \emph{étale} en ce point (18.4.9).
  \item[(iii)] Les notations et les hypothèses sur $Y$ étant celles de (18.10.1), supposons seulement que $f$ soit \emph{localement de type fini} (et non nécessairement localement de présentation finie). Alors, pour que $\sh{O}_{X,x}$ soit une $\sh{O}_{Y,y}$-algèbre \emph{essentiellement étale} (18.6.1), il faut et il suffit que $f$ vérifie les deux conditions suivantes:

\oldpage[IV-4]{159}

\begin{enumerate}
  \item[$1^\circ$] \itshape $f$ est formellement non ramifié au point $x$;
  \item[$2^\circ$] l'homomorphisme $\sh{O}_{Y,y}\to\sh{O}_{X,x}$ est injectif.
\end{enumerate}

Il n'y a encore à prouver que la suffisance de ces conditions. Tenant compte de (18.4.12), il suffit de montrer que $f$ est \emph{plat} au point $x$. Or, en reprenant la démonstration de (18.10.1) et les notations utilisées dans cette démonstration, il suffit de montrer que $f'$ est plat au point $x'$ (2.5.1); par ailleurs $f'$ est formellement non ramifié au point $x'$ (17.1.3); on peut donc se borner à faire la démonstration lorsque $Y=\Spec(A)$, où $A$ est \emph{strictement local} et $y=f(x)$ le point fermé de $Y$. Utilisant alors (18.8.3) on voit que l'homomorphisme $\sh{O}_{Y,y}\to\sh{O}_{X,x}$ est surjectif, donc \emph{bijectif} en vertu de l'hypothèse, et cela suffit à montrer la platitude de $f$ au point $x$.

Supposons de plus que $Y$ soit \emph{localement intègre} au point $y$ (\textbf{I}, 2.1.8). Alors les conditions $1^\circ$ et $2^\circ$ précédentes (jointe au fait que $Y$ est géométriquement unibranche au point $y$ et $f$ localement de type fini) entraînent déjà que $f$ est \emph{étale au point $x$}. En effet, cela résulte de ce qui précède et de (18.4.13).
\end{enumerate}
\end{remarks}

\begin{corollary}[18.10.3]
\label{IV.18.10.3}
Soient $Y$ un préschéma intègre et géométriquement unibranche, $\eta$ son point générique, $X$ un préschéma connexe, $f:X\to Y$ un morphisme localement de type fini et tel que $f^{-1}(\eta)$ soit non vide. Alors, si $f$ est formellement non ramifié, $f$ est étale, et $X$ est intègre et géométriquement unibranche.
\end{corollary}

Il résulte de l'hypothèse et de (18.4.13) que $f$ est étale en tous les points où il est \emph{plat}, et en particulier aux points de $f^{-1}(\eta)$, puisque $\sh{O}_{Y,\eta}=\boldsymbol{k}(\eta)$ est un corps. L'ensemble ouvert $U$ des points de $X$ où $f$ est étale est donc \emph{non vide}. Montrons d'autre part que pour tout $x\in\overline U$, l'homomorphisme $\sh{O}_{Y,f(x)}\to\sh{O}_{X,x}$ est \emph{injectif}. Soit en effet $s$ une section de $\sh{O}_Y$ au-dessus d'un voisinage ouvert de $f(x)$, qu'on peut toujours supposer être $Y$ lui-même (la question étant locale sur $X$ et $Y$), et supposons que son image $t\in\Gamma(Y,f_*(\sh{O}_X))=\Gamma(X,\sh{O}_X)$ ait un germe nul au point $x$, donc s'annule dans un voisinage ouvert $V$ de $x$; alors la restriction de $t$ à l'ouvert $U\cap V\neq\varnothing$ est nulle; mais la restriction $f|(U\cap V)$ est étale, donc $W=f(U\cap V)$ est ouvert dans $Y$ et contient par suite le point générique $\eta$; l'homomorphisme $\sh{O}_{Y,\eta}\to\sh{O}_{X,z}$ étant injectif pour tout point $z\in f^{-1}(\eta)$, l'hypothèse que $t|(U\cap V)=0$ entraîne que $s_\eta=0$, et comme $Y$ est intègre, cela entraîne $s=0$, d'où notre assertion. Appliquant maintenant (18.10.2, (iii)), on conclut que $f$ est \emph{plat} au point $x$, autrement dit $x\in U$, et $U$ est donc à la fois ouvert et fermé dans $X$; comme il est non vide et que $X$ est connexe, on a $U=X$.

Notons maintenant que puisque $f$ est étale, donc localement quasi-fini et que $Y$ est irréductible, les points maximaux de $X$ appartiennent à $f^{-1}(\eta)$ (2.3.4) et pour tout $x\in X$, il y a un voisinage de $x$ qui ne contient qu'un nombre fini de ces points maximaux; autrement dit l'ensemble des composantes irréductibles de $X$ est localement fini. Mais par (18.10.1) $X$ est intègre et géométriquement unibranche en tout point, donc tout point de $X$ n'appartient qu'à une seule composante irréductible, et puisque ces dernières forment un ensemble localement fini, elles sont ouvertes (et fermées) dans $X$; puisque $X$ est connexe, il est intègre.

